\begin{afterword}
\summaryhead

The post starts from the claim that human grouping instincts fit hunter-gatherer bands of about 50 people. That makes it surprising that large institutions exist at all, and the author names the kinds that do: governments that tax, corporations with central control of money, and religions spread by virulent memes. Most large institutions that would serve people's common interests, the post argues, ``fail to exist in the first place.''

Science is the example. Individuals rarely fund it, in the author's experience, except perhaps for diseases with gruesome victims. Although people are prosocial and give to puppy pounds, a science project is a poor emotional fit for personal giving: it takes years, its news is not stirring, only specialists can work on it, and donors see no immediate result. It seems to the author that very few things suit individual giving, and that this is why 200 million adult Americans find it hard to supervise the 535 members of Congress. Science is funded instead by governments, corporations and foundations, and scientists compete for the money. The post concludes that science survives as a ``parasite'' on the few large organizations that can exist, and that coordination beyond 50 people is ``too non-ancestral a problem.''

\responsehead

The post's thesis has a long history in economics, and that history supports it. Mancur Olson argued in 1965 that large groups do not act in their common interest without coercion or some other special device, and that large lobbies survive as by-products of organizations built for other purposes. The post's list of surviving institutions, and its picture of science attached to governments and corporations, match that account closely. Its explanation differs. Olson's people are self-interested and let others pay; the post grants that people are prosocial, and blames the poor fit between such tasks and the feelings that drive personal giving.

The evidence is one case, science funding, told from the author's fundraising experience and two examples of underfunded research. The case fits the post's account; one of its examples, the polywell device, was being funded by the US Navy at the time. The exception the post allows, research on diseases with ``gruesome victims,'' can be large: the March of Dimes funded the polio vaccine with small gifts. Its claim about foundations, that they spend money given for donors' reputations on science that ``promises to sound very charitable,'' is asserted without evidence.

The thesis in its general form cannot be checked directly. Institutions that ``fail to exist in the first place'' cannot be counted, so the claim that they are the ``vast majority'' is based on the one case.

In short: an account of why science depends on governments and large funders, in line with Olson's older theory of collective action, argued from one case.
\end{afterword}
